\chapter{Il quadro normativo e tecnico di riferimento}

\label{chap:quadro-normativo}

Per comprendere le scelte progettuali e implementative affrontate nei capitoli successivi, è necessario delineare il quadro di riferimento entro cui si colloca la verifica dell'accessibilità web.

Il capitolo introduce il concetto di accessibilità, i quattro principi delle WCAG e il rapporto tra gli standard tecnici e il quadro normativo europeo e italiano. Vengono inoltre illustrati il livello di conformità AA e i criteri WCAG considerati come riferimento per il lavoro.

\section{Cos'è l'accessibilità e qual è la sua finalità}

\label{sec:cose-accessibilita}

L'accessibilità digitale può essere definita, in primo luogo, richiamando la definizione fornita dalla Legge 9 gennaio 2004, n.~4, comunemente nota come legge Stanca: \enquote{la capacità dei sistemi informatici ivi inclusi i siti web e le applicazioni mobili, nelle forme e nei limiti consentiti dalle conoscenze tecnologiche, di erogare servizi e fornire informazioni fruibili, senza discriminazioni, anche da parte di coloro che a causa di disabilità necessitano di tecnologie assistive o configurazioni particolari}\legfootcite[art.~2, comma~1, lett.~a)]{legge:stanca}. La stessa legge collega esplicitamente tale obiettivo al principio di uguaglianza sancito dall'art.~3 della Costituzione, stabilendo che l'accesso delle persone con disabilità ai servizi informatici sia tutelato proprio \enquote{in ottemperanza} a tale principio\legfootcite[art.~1, comma~2]{legge:stanca}.

Una definizione più generale, di matrice tecnica anziché giuridica, è quella adottata dalla norma UNI CEI EN~301549:2021, che riprende a sua volta la norma ISO~9241-11, in cui l'accessibilità è definita come la \enquote{misura in cui i prodotti, i sistemi, i servizi, gli ambienti e le strutture possono essere utilizzati da persone di una popolazione con la più ampia gamma di necessità degli utenti, caratteristiche e capacità, per raggiungere un obiettivo identificato in un contesto identificato di utilizzo}\footcite[cap.~3, voce \enquote{accessibilità}, da ISO~9241-11:2018]{norma:en301549}. Le due definizioni si completano: la prima pone in evidenza la necessità di garantire l'accesso alle persone con disabilità, anche attraverso tecnologie assistive o configurazioni particolari, e colloca tale obiettivo nel quadro del principio di uguaglianza; la seconda descrive l'accessibilità in termini più ampi, come proprietà di un sistema rispetto a una popolazione di utenti eterogenea per caratteristiche e capacità, di cui le persone con disabilità sono una componente e non la totalità.

Questa prospettiva più ampia emerge con particolare chiarezza nella struttura della norma EN~301549, in particolare nel capitolo~4, dedicato alle prestazioni funzionali che un sistema accessibile deve garantire\footcite[cap.~4 \enquote{Prestazioni funzionali}]{norma:en301549}. Per ciascuna delle dieci prestazioni individuate, la norma specifica la categoria di utenti per cui il requisito risulta essenziale e, al contempo, evidenzia come la medesima caratteristica possa apportare benefici anche in altre situazioni. Ad esempio, per l'utilizzo senza vista, la norma richiede che l'ICT fornisca una modalità operativa che non richieda la vista e precisa: \enquote{Questo è essenziale per gli utenti senza vista e avvantaggia molti altri utenti in diverse situazioni}\footcite[§~4.2.1 \enquote{Utilizzo senza vista}]{norma:en301549}. Una formulazione pressoché identica ricorre nelle successive sezioni, ciascuna dedicata a una diversa condizione di abilità, adattando di volta in volta il riferimento alla specifica condizione considerata.

Il principio è facilmente osservabile anche in caratteristiche di uso comune: i sottotitoli sono essenziali per chi presenta una perdita o una riduzione della capacità uditiva, ma possono essere utili anche a chi guarda un video in un ambiente rumoroso o senza audio disponibile; un'alternativa testuale a un'immagine è fondamentale per chi utilizza uno screen reader, ma può risultare utile anche quando l'immagine non viene caricata correttamente; un adeguato contrasto tra testo e sfondo è particolarmente importante per alcune esigenze visive, ma migliora la leggibilità anche in condizioni ambientali sfavorevoli. L'accessibilità non consiste quindi soltanto nel rimuovere ostacoli per una specifica categoria di utenti, ma nel progettare sistemi utilizzabili in una maggiore varietà di condizioni.

Questa impostazione trova riscontro anche nella normativa europea: la Direttiva (UE) 2019/882, relativa ai requisiti di accessibilità dei prodotti e dei servizi, definisce le persone con disabilità in coerenza con la Convenzione delle Nazioni Unite sui diritti delle persone con disabilità, riconoscendo che l'interazione tra caratteristiche individuali e barriere di varia natura può ostacolare la piena partecipazione alla società\legfootcite[considerando~3]{direttiva:eaa}. La stessa direttiva riconosce inoltre che i requisiti di accessibilità possono apportare benefici anche ad altre persone con limitazioni funzionali, tra cui persone anziane, donne in gravidanza e persone che viaggiano con bagagli, considerando anche limitazioni permanenti o temporanee\legfootcite[considerando~4]{direttiva:eaa}.

La necessità di considerare condizioni d'uso differenti porta quindi a superare una concezione dell'accessibilità legata esclusivamente a condizioni permanenti. Una barriera può infatti essere associata a una condizione permanente, come una difficoltà visiva; a una condizione temporanea, come un braccio ingessato che impedisce l'utilizzo del mouse; oppure a una condizione situazionale, come la difficoltà di leggere uno schermo esposto alla luce diretta del sole, di ascoltare un contenuto audio in un ambiente rumoroso o di utilizzare le mani mentre si tiene in braccio un bambino. In tutti questi casi, la barriera può emergere dall'interazione tra le caratteristiche dell'utente, il contesto e le modalità con cui il sistema è stato progettato.

Questa prospettiva costituisce uno dei riferimenti concettuali che guidano il presente lavoro. Un'estensione finalizzata alla verifica dell'accessibilità di una pagina web risulta tanto più significativa quanto più i criteri selezionati sono in grado di intercettare barriere che possono manifestarsi in differenti condizioni d'uso. L'obiettivo dell'accessibilità è quindi contribuire alla realizzazione di interfacce che possano essere utilizzate da una platea più ampia e in una maggiore varietà di situazioni.

\section{Cosa deve garantire il web: i quattro principi dell'accessibilità}

\label{sec:quattro-principi}

Considerando che l'accessibilità riguarda una platea di utenti così eterogenea, è necessario disporre di un quadro di riferimento che permetta di valutare in modo sistematico la fruibilità di un contenuto digitale. Questo quadro è fornito dalle \gls{wcag} del World Wide Web Consortium (W3C), le linee guida di riferimento per l'accessibilità dei contenuti web, organizzate attorno a quattro principi fondamentali, noti con l'acronimo POUR: Percepibile, Operabile, Comprensibile e Robusto\footcite[§ \enquote{Understanding the Four Principles of Accessibility}]{w3c:understanding-intro}.

Il documento ufficiale che ne illustra il significato definisce ciascun principio attraverso una formulazione sintetica, cui associa una spiegazione delle condizioni necessarie per rendere il contenuto accessibile\footcite[§ \enquote{Understanding the Four Principles of Accessibility}]{w3c:understanding-intro}. Un contenuto è \emph{percepibile} quando \enquote{le informazioni e i componenti dell'interfaccia utente devono essere presentabili agli utenti in modi che essi possano percepire}: l'utente deve quindi poter accedere alle informazioni attraverso almeno una modalità sensoriale compatibile con le proprie capacità. Un contenuto è \emph{operabile} quando \enquote{i componenti dell'interfaccia utente e la navigazione [sono] operabili}: l'interazione non può quindi dipendere esclusivamente da una modalità di input che l'utente non è in grado di utilizzare. Un contenuto è \emph{comprensibile} quando \enquote{le informazioni e il funzionamento dell'interfaccia utente [sono] comprensibili}: non è sufficiente che le informazioni siano accessibili, ma anche il loro significato e il funzionamento dell'interfaccia devono poter essere compresi. Infine, un contenuto è \emph{robusto} quando \enquote{il contenuto [è] sufficientemente solido da poter essere interpretato in modo affidabile da un'ampia varietà di programmi utente, incluse le tecnologie assistive}: il contenuto deve quindi mantenere una corretta interpretabilità attraverso differenti tecnologie, comprese quelle assistive\footcite{w3c:understanding-intro}. I quattro principi costituiscono così la base sulla quale vengono costruiti i requisiti più specifici delle WCAG.

Nella struttura delle WCAG, dai quattro principi discendono tredici linee guida (\emph{guidelines}), che definiscono obiettivi generali, e, per ciascuna di esse, uno o più criteri di successo (\emph{success criteria}), formulati invece come requisiti verificabili e applicabili a contenuti specifici, indipendentemente dalla tecnologia utilizzata\footcite[§~0.2 \enquote{WCAG 2 Livelli di orientamento}]{w3c:wcag21}. Sono i criteri di successo a determinare la conformità alle WCAG e a essere associati ai tre livelli di conformità A, AA e AAA, discussi nella sezione~\ref{sec:livello-aa}. A questa struttura si affiancano le tecniche sufficienti e consigliate, documentate separatamente dal W3C come materiale di supporto per l'implementazione e la verifica dei criteri\footcite{w3c:quickref}.

I quattro principi, pertanto, non costituiscono soltanto il modello concettuale attraverso cui il W3C organizza le WCAG, ma forniscono un riferimento che ricorre anche nel quadro normativo considerato nel presente lavoro. La legge Stanca, nella versione aggiornata a seguito del recepimento della direttiva europea sull'accessibilità dei siti web e delle applicazioni mobili degli enti pubblici, stabilisce che i siti web e le applicazioni mobili dei soggetti erogatori siano accessibili se risultano \enquote{percepibili, utilizzabili, comprensibili e solidi}\legfootcite[art.~3-bis, comma~1]{legge:stanca}.

La norma tecnica europea UNI CEI EN~301549:2021, nella versione italiana del febbraio 2022, adotta invece una prospettiva complementare, descrivendo nel capitolo~4 le prestazioni funzionali associate alle esigenze di differenti utenti. In particolare, il punto~4.2 descrive dieci prestazioni funzionali relative a diverse modalità di utilizzo dell'ICT, alle quali si aggiunge una prestazione relativa alla privacy\footcite[cap.~4, §~4.2]{norma:en301549}. Tali prestazioni descrivono le esigenze funzionali degli utenti e sono collegate ai requisiti tecnici previsti dalle altre parti della norma\footcite[cap.~4]{norma:en301549}.

Il Decreto Legislativo 27 maggio 2022, n.~82, che recepisce in Italia la Direttiva (UE) 2019/882 sui requisiti di accessibilità dei prodotti e dei servizi, contiene nell'Allegato~I, Sezione~VII, i criteri funzionali di prestazione, applicabili nei casi e secondo le condizioni previste dalla stessa sezione\legfootcite[Allegato~I, Sezione~VII]{dlgs:82-2022}. Le più recenti linee guida AGID, nel presentare i \enquote{quattro pilastri dell'accessibilità}, richiamano invece esplicitamente i quattro principi POUR\footcite[§~6.1]{agid:lg-82-2022}. I diversi riferimenti utilizzano quindi prospettive complementari: i principi POUR forniscono una struttura generale per descrivere l'accessibilità dei contenuti web, mentre le prestazioni funzionali descrivono esigenze e modalità d'uso cui devono rispondere i prodotti e i servizi ICT.

La continuità tra questi diversi livelli di riferimento mostra come i quattro principi forniscano una chiave di lettura comune delle barriere digitali: la difficoltà può riguardare la possibilità di percepire le informazioni, di interagire con l'interfaccia, di comprenderne il contenuto o il funzionamento, oppure di interpretare correttamente il contenuto attraverso le tecnologie utilizzate. Non si tratta quindi di una semplice suddivisione teorica, ma di un modello che consente di ricondurre esigenze e requisiti differenti a quattro dimensioni fondamentali dell'accessibilità.

È proprio questa articolazione a rendere i quattro principi rilevanti ai fini della selezione dei criteri \gls{wcag} oggetto di verifica automatica nel presente lavoro. Come si vedrà nella sezione~\ref{sec:selezione-criteri}, la selezione è stata condotta cercando di rappresentare, per quanto possibile, tutti e quattro i principi, considerando l'accessibilità come una proprietà multidimensionale nella quale percepibilità, operabilità, comprensibilità e robustezza concorrono alla fruibilità del contenuto.

\section{Come questi principi diventano obbligo giuridico: la stratificazione delle fonti}

\label{sec:stratificazione-fonti}

I principi dell'accessibilità web trovano applicazione all'interno di un sistema composto da fonti di natura e livello differenti. Le WCAG definiscono il quadro tecnico per l'accessibilità dei contenuti web; le norme tecniche europee collegano i requisiti di accessibilità al più ampio insieme dei prodotti e servizi ICT; le direttive europee stabiliscono obblighi e ambiti di applicazione; infine, la normativa italiana e le linee guida AGID ne disciplinano l'applicazione nell'ordinamento nazionale. Distinguere questi livelli è necessario per comprendere perché un requisito tecnico possa assumere rilevanza giuridica senza che le WCAG costituiscano, di per sé, una fonte legislativa.

\subsection{La legge Stanca: la prima normativa italiana}

L'Italia si è dotata già nel 2004 di una disciplina organica sull'accessibilità degli strumenti informatici. La legge 9 gennaio 2004, n.~4, nota come legge Stanca, riconosce il diritto di accesso alle fonti di informazione e ai relativi servizi attraverso gli strumenti informatici e telematici, con particolare riferimento alle persone con disabilità\legfootcite[art.~1]{legge:stanca}. La disciplina è stata successivamente modificata per adeguarsi all'evoluzione tecnologica e al quadro europeo, includendo espressamente siti web e applicazioni mobili e, in determinate condizioni, anche soggetti privati di grandi dimensioni\legfootcite[art.~3, commi~1 e 1-bis]{legge:stanca}.

Un elemento caratterizzante della legge è la separazione tra il principio giuridico dell'accessibilità e la sua specificazione tecnica. L'articolo~11 affida infatti ad AGID l'emanazione di linee guida nelle quali sono stabiliti i requisiti tecnici per l'accessibilità degli strumenti informatici, comprese le metodologie tecniche per la verifica\legfootcite[art.~11]{legge:stanca}. La legge definisce quindi il quadro dell'obbligo, mentre il dettaglio tecnico viene disciplinato attraverso linee guida aggiornabili in relazione all'evoluzione degli standard e delle tecnologie.

Le Linee guida sull'accessibilità degli strumenti informatici costituiscono il riferimento operativo per l'applicazione della disciplina. Il documento AGID richiama espressamente le WCAG~2.1 e la norma UNI CEI EN~301549:2021, nella versione italiana del febbraio 2022\footcite{agid:lg-stanca}. Per le pagine web, le linee guida stabiliscono che la conformità al livello AA delle WCAG~2.1 sia equivalente alla conformità ai requisiti relativi al web individuati nei punti da 9.1 a 9.4 della EN~301549 e ai relativi requisiti di conformità\footcite[§~2.2]{agid:lg-stanca}. Le WCAG assumono così un ruolo tecnico centrale nell'applicazione della disciplina, pur rimanendo una specifica elaborata dal W3C e non una fonte legislativa.

\subsection{Il livello europeo: dalla Direttiva 2016/2102 all'European Accessibility Act}

Il quadro europeo si è sviluppato attraverso due interventi distinti. Il primo è la Direttiva (UE) 2016/2102, relativa all'accessibilità dei siti web e delle applicazioni mobili degli enti pubblici\legfootcite{direttiva:2016-2102}, che stabilisce requisiti comuni per l'accessibilità dei contenuti digitali degli organismi del settore pubblico e collega tali requisiti al quadro tecnico europeo rappresentato dalla EN~301549\legfootcite[art.~6]{direttiva:2016-2102}.

La norma EN~301549\footcite{norma:en301549} è una norma tecnica europea armonizzata attraverso la quale i requisiti di accessibilità possono essere applicati e verificati in modo uniforme. Per il web, il punto~9 della norma stabilisce il collegamento con le WCAG~2.1: il punto~9.0 specifica che la conformità al livello AA delle WCAG~2.1 è equivalente alla conformità ai requisiti dei punti da 9.1 a 9.4 e ai requisiti di conformità previsti dal punto~9.6\footcite[§~9.0]{norma:en301549}. La norma costituisce quindi il principale raccordo tecnico tra i requisiti di accessibilità e i criteri di successo delle WCAG applicabili alle pagine web.

Il secondo intervento europeo è la Direttiva (UE) 2019/882, nota come European Accessibility Act (EAA), che introduce requisiti di accessibilità per specifiche categorie di prodotti e servizi destinati ai consumatori\legfootcite{direttiva:eaa}. Tra i servizi interessati rientrano, tra gli altri, i servizi di comunicazione elettronica, alcuni servizi di trasporto passeggeri, i servizi bancari per consumatori, i servizi di commercio elettronico e alcuni servizi relativi ai libri elettronici\legfootcite[art.~2]{direttiva:eaa}. L'EAA amplia quindi il quadro europeo oltre l'ambito specifico dei siti web e delle applicazioni mobili degli enti pubblici, introducendo requisiti per un insieme definito di prodotti e servizi immessi sul mercato o forniti ai consumatori.

\subsection{Il recepimento italiano: il D.Lgs.~82/2022 e l'aggiornamento della legge Stanca}

La Direttiva (UE) 2019/882 è stata recepita in Italia attraverso il Decreto Legislativo 27 maggio 2022, n.~82, che stabilisce i requisiti di accessibilità per specifiche categorie di prodotti e servizi rientranti nel proprio ambito di applicazione, con disposizioni applicabili a decorrere dal 28 giugno 2025\legfootcite[art.~1, comma~1]{dlgs:82-2022}. Il decreto individua, tra i prodotti interessati, sistemi hardware e sistemi operativi informatici generici per consumatori, terminali self-service, apparecchiature terminali utilizzate per servizi di comunicazione elettronica o per accedere a servizi di media audiovisivi e lettori di libri elettronici. Per quanto riguarda i servizi, rientrano nell'ambito di applicazione, tra gli altri, i servizi di comunicazione elettronica, alcuni servizi di trasporto passeggeri, i servizi bancari per consumatori, i servizi di commercio elettronico e alcuni servizi relativi ai libri elettronici\legfootcite[art.~1, commi~2-3]{dlgs:82-2022}.

Il criterio adottato dal D.Lgs.~82/2022 è quindi differente rispetto a quello della legge Stanca. Quest'ultima si applica alle pubbliche amministrazioni e agli altri soggetti pubblici individuati dall'articolo~3 e, alle condizioni previste dalla legge, anche ai soggetti privati con un fatturato medio superiore a 500 milioni di euro negli ultimi tre anni di attività\legfootcite[art.~3, commi~1 e 1-bis]{legge:stanca}. Il D.Lgs.~82/2022 individua invece il proprio ambito di applicazione sulla base di specifiche categorie di prodotti e servizi, cui si applicano i requisiti di accessibilità secondo le condizioni previste dal decreto. Le due discipline presentano quindi ambiti di applicazione differenti e non costituiscono un unico regime applicabile indistintamente a tutti i siti web.

Nel caso del D.Lgs.~82/2022, i requisiti di accessibilità sono definiti nell'Allegato~I, che raccoglie le prescrizioni applicabili ai prodotti e ai servizi compresi nell'ambito di applicazione del decreto\legfootcite[Allegato~I]{dlgs:82-2022}. La Sezione~VII dell'Allegato è dedicata ai criteri funzionali di prestazione e stabilisce che, quando le sezioni precedenti non disciplinano una o più funzioni del prodotto o del servizio, tali funzioni debbano essere accessibili attraverso il rispetto dei criteri funzionali pertinenti, secondo le condizioni previste dalla stessa sezione\legfootcite[Allegato~I, Sezione~VII]{dlgs:82-2022}. Il decreto definisce quindi i requisiti sul piano normativo, mentre le linee guida adottate in attuazione dell'articolo~21 ne supportano l'interpretazione e l'applicazione.

In questo quadro si inseriscono le Linee guida sull'accessibilità dei servizi adottate da AGID nel 2026 in attuazione dell'articolo~21 del D.Lgs.~82/2022. Le linee guida includono la EN~301549 tra le norme tecniche di riferimento e, nella sezione dedicata alla norma, ne illustrano il ruolo nell'accessibilità dei siti web e delle applicazioni mobili degli enti pubblici nell'ambito della Direttiva (UE) 2016/2102, evidenziandone inoltre i numerosi elementi comuni con le WCAG~2.1\footcite[§§~2.2 e 8.5]{agid:lg-82-2022}. Il riferimento alla EN~301549 si inserisce quindi nel più ampio quadro tecnico europeo dell'accessibilità, senza coincidere con l'individuazione dell'ambito di applicazione del D.Lgs.~82/2022.

La legge Stanca, parallelamente, è rimasta in vigore ed è stata modificata nel tempo per adeguarsi all'evoluzione del quadro europeo e tecnico. In particolare, il recepimento della Direttiva (UE) 2016/2102 ha portato all'inserimento nella legge di disposizioni specifiche relative all'accessibilità dei siti web e delle applicazioni mobili degli enti pubblici. L'attuale articolo~11 continua ad affidare ad AGID l'emanazione delle linee guida nelle quali sono stabiliti i requisiti tecnici e le metodologie per la verifica dell'accessibilità degli strumenti informatici, inclusi siti web e applicazioni mobili\footcite[art.~11]{legge:stanca}.

Si configurano quindi due percorsi normativi distinti: da un lato, la disciplina della legge Stanca, sviluppata anche attraverso il recepimento della Direttiva (UE) 2016/2102 e rivolta ai soggetti rientranti nel relativo ambito di applicazione; dall'altro, la disciplina introdotta dal D.Lgs.~82/2022 in attuazione dell'European Accessibility Act, applicabile alle specifiche categorie di prodotti e servizi individuate dal decreto. Pur avendo ambiti di applicazione differenti, entrambe si inseriscono in un quadro tecnico europeo nel quale la norma EN~301549 e le WCAG costituiscono riferimenti rilevanti per la definizione e la verifica dei requisiti di accessibilità.

\section{Dal principio tecnico al requisito verificabile}

La stratificazione delle fonti può essere quindi ricondotta a una sequenza nella quale ogni livello svolge una funzione diversa. Le WCAG definiscono i principi, le linee guida e i criteri di successo attraverso i quali descrivere l'accessibilità dei contenuti web; la norma EN~301549 collega questi requisiti al più ampio quadro dei prodotti e servizi ICT e ne definisce l'applicazione nei diversi ambiti; le direttive europee stabiliscono gli obblighi e i relativi ambiti di applicazione; infine, la normativa italiana e le linee guida AGID ne disciplinano l'applicazione nel contesto nazionale.

Questa distinzione è importante per evitare di confondere il ruolo delle diverse fonti. Le WCAG non stabiliscono autonomamente quali soggetti siano giuridicamente obbligati all'accessibilità. Analogamente, la EN~301549 non individua da sola i destinatari degli obblighi: fornisce piuttosto un riferimento tecnico attraverso il quale i requisiti possono essere specificati e verificati. Sono le disposizioni normative a determinare quando tali requisiti assumono rilevanza giuridica per un determinato soggetto, prodotto o servizio.

Per il presente lavoro, il passaggio più rilevante è quindi quello che collega il livello normativo al livello della verifica automatica. Una volta individuato nel quadro tecnico-normativo il riferimento alle WCAG~2.1 livello AA per le pagine web, i criteri di successo delle WCAG possono essere utilizzati come unità di riferimento per individuare condizioni di accessibilità concretamente verificabili all'interno di una pagina. La selezione dei criteri analizzata nei capitoli successivi si colloca precisamente a questo livello: non mira a stabilire autonomamente la conformità giuridica di un sito, ma a individuare quali requisiti delle WCAG possano essere tradotti in controlli utili alla valutazione dell'accessibilità.

\section{Perché il livello AA e perché solo alcuni dei criteri WCAG}

\label{sec:livello-aa}

Le \gls{wcag} articolano i criteri di successo su tre livelli di conformità crescente: A, AA e AAA. I livelli non rappresentano una graduatoria di importanza dei criteri, ma definiscono requisiti cumulativi di conformità: il livello AA comprende tutti i criteri di livello A e AA, mentre il livello AAA comprende anche quelli di livello AAA\footcite[§~5.2.1]{w3c:wcag21}. Il W3C precisa inoltre che non è raccomandato richiedere il livello AAA come requisito generale per interi siti, poiché non è possibile soddisfare tutti i relativi criteri per alcune tipologie di contenuto\footcite[§~5.2.1]{w3c:wcag21}.

Nel contesto normativo considerato in questo lavoro, il riferimento tecnico per le pagine web è il livello AA delle WCAG~2.1. Le Linee guida AGID sull'accessibilità degli strumenti informatici stabiliscono infatti che la conformità al livello AA delle WCAG~2.1 sia equivalente alla conformità ai requisiti della EN~301549 relativi alle pagine web\footcite[§~2.2]{agid:lg-stanca}. La stessa impostazione è riportata nelle linee guida AGID dedicate ai soggetti privati\footcite{agid:lg-stanca-1bis}. Per questo motivo, il presente lavoro assume le WCAG~2.1 livello AA come riferimento tecnico per la selezione dei criteri da sottoporre a verifica.

Per le pagine web, i criteri di successo di livello A e AA delle WCAG~2.1 costituiscono complessivamente un insieme di 50 criteri. Il dato è confermato anche da AGID, che nel descrivere il proprio monitoraggio approfondito delle pagine web indica esplicitamente la verifica di tutti i 50 criteri di successo di livello A e AA\footcite{agid:lg-stanca}. Questo insieme rappresenta il riferimento completo dal quale il presente lavoro seleziona un sottoinsieme di criteri adatti alla verifica automatica assistita.

La scelta di non verificare la totalità dei criteri nasce da una caratteristica intrinseca dell'accessibilità web: non tutti i requisiti sono ugualmente verificabili attraverso procedure automatiche. Alcuni criteri possono essere controllati direttamente attraverso l'analisi della struttura e del codice della pagina; altri richiedono invece una valutazione del significato o dell'adeguatezza del contenuto; altri ancora dipendono dal comportamento dell'interfaccia durante l'interazione e richiedono quindi una valutazione dinamica. Un sistema automatico può pertanto fornire un supporto significativo alla valutazione dell'accessibilità, ma non sostituisce necessariamente tutte le forme di verifica previste dalle WCAG.

La selezione dei criteri da includere nell'estensione è stata quindi condotta incrociando due esigenze. Da un lato, si è considerata la frequenza con cui determinate problematiche di accessibilità emergono sul web reale, facendo riferimento ai dati del WebAIM Million\footcite{webaim:million2026} e dei monitoraggi AGID\footcite{agid:errori-monitoraggio}; dall'altro, è stata valutata la possibilità di tradurre il relativo requisito in un controllo compatibile con un'architettura basata su \gls{llm}. Il primo elemento permette di concentrare l'attenzione su problematiche effettivamente diffuse, mentre il secondo consente di mantenere la selezione coerente con l'obiettivo tecnico del progetto.

A questi due criteri si aggiunge la prospettiva già introdotta nella sezione~\ref{sec:cose-accessibilita}: a parità di frequenza e fattibilità tecnica, risultano particolarmente interessanti i criteri in grado di intercettare barriere che possono manifestarsi in una pluralità di condizioni d'uso. La selezione non mira quindi semplicemente a individuare i controlli più semplici da automatizzare, ma a costruire un sottoinsieme che rappresenti, per quanto possibile, le diverse dimensioni dell'accessibilità e che sia coerente con l'obiettivo pratico dell'estensione.

I criteri specifici selezionati e la metodologia con cui sono stati individuati a partire dai dati del WebAIM Million e dei monitoraggi AGID sono descritti nella sezione~\ref{sec:selezione-criteri}.

\section{L'accessibilità come metodo di progettazione, non come rifinitura finale}

\label{sec:accessibilita-by-design}

Quanto ricostruito nelle sezioni precedenti permette di formulare un'ultima considerazione di carattere metodologico. L'accessibilità, così come emerge dalle definizioni esaminate nella sezione~\ref{sec:cose-accessibilita} e dalla logica dei quattro principi POUR discussa nella sezione~\ref{sec:quattro-principi}, non dovrebbe essere considerata come una caratteristica da aggiungere a un prodotto digitale già completato, ma come una proprietà da considerare durante l'intero processo di progettazione e sviluppo.

Le Linee guida AGID distinguono infatti un approccio \enquote{by design}, nel quale l'accessibilità viene integrata fin dalle fasi iniziali del ciclo di vita del prodotto o servizio, da un approccio correttivo che interviene successivamente\footcite[§~6.2]{agid:lg-82-2022}. Questa impostazione è coerente con la struttura stessa dei requisiti di accessibilità: intervenire durante la progettazione consente di considerare fin dall'inizio differenti modalità di percezione, interazione e comprensione, mentre la verifica a posteriori permette di individuare e correggere barriere già presenti.

La necessità di entrambe le prospettive emerge anche osservando lo stato del web reale. Il WebAIM Million analizza annualmente le home page di un milione di siti selezionati sulla base del ranking Tranco mediante il motore WAVE. L'analisi non costituisce una verifica completa della conformità alle WCAG, poiché considera gli errori rilevabili automaticamente, ma fornisce un'indicazione quantitativa sulla diffusione di alcune barriere osservabili attraverso strumenti automatici\footcite{webaim:million2026}. Nell'edizione 2026, relativa alle analisi condotte nel febbraio dello stesso anno, il 95,9% delle home page analizzate presentava almeno un errore WCAG rilevabile automaticamente, con una media di 56,1 errori rilevati per pagina\footcite{webaim:million2026}. Lo stesso report sottolinea che l'assenza di errori rilevati automaticamente non implica che una pagina sia complessivamente accessibile o conforme alle WCAG.

Il dato evidenzia quindi una distanza tra la disponibilità di criteri e strumenti per la progettazione accessibile e la loro applicazione concreta sul web. La verifica automatica non può colmare da sola questa distanza, ma può contribuire a rendere più sistematico il controllo di una parte delle barriere rilevabili attraverso procedure automatiche. In questo senso, l'automazione non sostituisce la valutazione complessiva dell'accessibilità, ma può costituire uno strumento di supporto per individuare problematiche ricorrenti e indirizzare le successive attività di verifica e correzione.

È in questo secondo fronte, quello della verifica di contenuti web già esistenti, che si colloca il problema pratico affrontato nel presente lavoro. L'obiettivo dell'estensione non è quindi sostituire una valutazione completa dell'accessibilità, ma fornire un ulteriore livello di supporto alla verifica, concentrato su un sottoinsieme di criteri WCAG selezionati sulla base della loro rilevanza e della possibilità di automatizzarne il controllo.